\documentclass[10pt]{article}
\usepackage[margin=0.88in]{geometry}
\usepackage[T1]{fontenc}
\usepackage{lmodern}
\usepackage{amsmath,amssymb,amsthm,graphicx,microtype}
\usepackage[hidelinks]{hyperref}

\newtheorem{theorem}{Theorem}
\newtheorem{lemma}[theorem]{Lemma}
\newcommand{\N}{\mathbb N}
\newcommand{\I}{\mathcal I}
\newcommand{\J}{\mathcal J}
\newcommand{\M}{\mathcal M}
\newcommand{\U}{\mathcal U}
\newcommand{\Fin}{\mathsf{Fin}}
\newcommand{\linf}{\ell_\infty}

\title{A Complemented Pair of Ideal-Null Spaces\\with Infinite-Dimensional Quotient}
\author{Candidate full solution packet\\
\small arXiv:2507.13866, Remark 6.7}
\date{11 August 2026}

\begin{document}
\maketitle

\noindent\textbf{Status.} Candidate full affirmative answer, subject to
expert review.  We construct proper ideals
$\Fin\subsetneq\I\subsetneq\J$ on $\N$ such that $c_{0,\I}$ is
complemented in $c_{0,\J}$ and, more strongly,
\[
 c_{0,\J}/c_{0,\I}\quad\hbox{is linearly isometric to }\linf.
\]

\section{Source question}

For an ideal $\mathcal A$ on $\N$, containing every finite set, put
\[
 c_{0,\mathcal A}
 =\left\{x\in\linf:
   \{k:|x(k)|\geq\varepsilon\}\in\mathcal A
   \text{ for every }\varepsilon>0\right\}.
\]
Rinc\'on-Villamizar and Uzc\'ategui Aylwin study when
$c_{0,\I}$ is complemented in $c_{0,\J}$ for proper ideals
$\I\subsetneq\J$ \cite{RU}.  After constructing examples with every finite
quotient dimension, they ask in Remark 6.7 on printed page 17:

\begin{center}
\includegraphics[width=0.98\textwidth]{figures/open_problem_crop.png}
\end{center}

\begin{quote}
``We do not know if there exist examples with
$\dim(c_{0,\mathcal J}/c_{0,\mathcal I})=\infty$.''
\end{quote}

\section{Proof intuition}

A maximal ideal is the complement of an ultrafilter, and its ideal-null
space is the kernel of the corresponding ultralimit functional.  Put
countably many free ultrafilters on pairwise disjoint infinite blocks of
$\N$.  The map collecting their ultralimits is a contraction from
$\linf$ to $\linf$.  It has an explicit isometric right inverse: make a
sequence constant on each block.

Take $\I$ to be the intersection of all the associated maximal ideals and
$\J$ to be one distinguished maximal ideal.  On $c_{0,\J}$, the remaining
ultralimits still map onto $\linf$, their joint kernel is exactly
$c_{0,\I}$, and the block-constant right inverse yields a projection of norm
at most two.  The same right inverse identifies the quotient norm exactly,
so the quotient is not merely infinite-dimensional but isometric to
$\linf$.

\section{Construction and proof}

Partition $\N$ into pairwise disjoint infinite sets
\[
 \N=\bigsqcup_{n=0}^{\infty}A_n.
\]
For each $n$, choose a free ultrafilter $\mathcal V_n$ on $A_n$ and lift it
to the ultrafilter $\U_n$ on $\N$ defined by
\[
 E\in\U_n\quad\Longleftrightarrow\quad E\cap A_n\in\mathcal V_n.
\]
Let $\M_n=\mathcal P(\N)\setminus\U_n$ be the corresponding maximal ideal,
and define
\[
 \J=\M_0,
 \qquad
 \I=\bigcap_{n=0}^{\infty}\M_n.
 \tag{1}\label{eq:ideals}
\]

\begin{lemma}\label{lem:kernels}
For $x\in\linf$, let
\[
 \phi_n(x)=\lim_{\U_n}x.
\]
Then $\phi_n$ is a norm-one linear functional and
\[
 c_{0,\M_n}=\ker\phi_n,
 \qquad
 c_{0,\I}=\bigcap_{n=0}^{\infty}\ker\phi_n,
 \qquad
 c_{0,\J}=\ker\phi_0.
 \tag{2}\label{eq:kernels}
\]
\end{lemma}

\begin{proof}
The existence and linearity of the ultralimit of a bounded scalar sequence
are standard, and $\phi_n(\mathbf 1)=1$ gives norm one.  For
$\varepsilon>0$, the assertion $\phi_n(x)=0$ is equivalent to
\[
 \{k:|x(k)|<\varepsilon\}\in\U_n,
\]
for every $\varepsilon>0$.  Taking complements, this is equivalent to
$\{k:|x(k)|\geq\varepsilon\}\in\M_n$ for every $\varepsilon$, which is
precisely $x\in c_{0,\M_n}$.  Intersecting these equivalences and using
\eqref{eq:ideals} proves the other two identities.
\end{proof}

\begin{theorem}\label{thm:main}
The ideals in \eqref{eq:ideals} are proper and satisfy
\[
 \Fin\subsetneq\I\subsetneq\J.
\]
There is a projection
\[
 P:c_{0,\J}\longrightarrow c_{0,\I}
 \quad\text{with}\quad \|P\|\leq2,
\]
and $c_{0,\J}/c_{0,\I}$ is linearly isometric to $\linf$.
Consequently the question in Remark 6.7 of \cite{RU} has an affirmative
answer.
\end{theorem}

\begin{proof}
Each $\M_n$ contains $\Fin$, since $\U_n$ is free.  Choose one point
$s_n\in A_n$ for each $n$ and let $S=\{s_n:n\geq0\}$.  For every $m$,
$S\cap A_m$ is a singleton and hence is not in the free ultrafilter
$\mathcal V_m$.  Thus $S\in\M_m$ for all $m$, so the infinite set $S$
belongs to $\I$ and $\Fin\subsetneq\I$.  Moreover, $A_1\in\M_0=\J$ but
$A_1\notin\M_1$, giving $\I\subsetneq\J$.  Both ideals are proper because
$\N$ belongs to every $\U_n$.

Define
\[
 L:c_{0,\J}=\ker\phi_0\longrightarrow\linf(\{1,2,\ldots\}),
 \qquad
 Lx=(\phi_n(x))_{n\geq1}.
 \tag{3}\label{eq:L}
\]
This is a contraction.  For $y=(y_n)_{n\geq1}\in\linf$, define the
block-constant sequence $Ry\in\linf(\N)$ by
\[
 (Ry)(k)=
 \begin{cases}
 0,&k\in A_0,\\
 y_n,&k\in A_n,\ n\geq1.
 \end{cases}
 \tag{4}\label{eq:R}
\]
Then $\|Ry\|=\|y\|$, $\phi_0(Ry)=0$, and
$\phi_n(Ry)=y_n$ for $n\geq1$.  Therefore $R$ is an isometric right inverse
of $L$ whose range lies in $c_{0,\J}$.

By Lemma~\ref{lem:kernels},
\[
 \ker L=\bigcap_{n=0}^{\infty}\ker\phi_n=c_{0,\I}.
\]
It follows that
\[
 P=I_{c_{0,\J}}-RL
 \tag{5}\label{eq:P}
\]
is a projection onto $c_{0,\I}$ and $\|P\|\leq2$.

Finally, let $\widetilde L:c_{0,\J}/c_{0,\I}\to\linf$ be the map induced
by $L$.  If $x\in c_{0,\J}$ and $z\in c_{0,\I}$, then
\[
 \|Lx\|=\|L(x+z)\|\leq\|x+z\|,
\]
so $\|Lx\|\leq\|x+c_{0,\I}\|$.  Conversely,
$x-RLx=Px\in c_{0,\I}$, and hence
\[
 \|x+c_{0,\I}\|
 \leq\|RLx\|=\|Lx\|.
\]
Thus $\widetilde L$ is an onto linear isometry, completing the proof.
\end{proof}

\section{Verification and limitations}

The construction uses free ultrafilters, hence the usual ultrafilter lemma
available in ZFC.  It gives an explicit relative projection of norm at most
two and identifies the quotient isometrically; no separability, positivity
of the relative projection, or definability of the ideals is claimed.
Indeed the ideals are necessarily nonmeager and highly nondefinable, in line
with the obstructions developed in the source.

The main verifier checks are: (i) the level-set definition of
$c_{0,\M_n}$ agrees with the kernel of the scalar ultralimit; (ii) the
block-constant right inverse lies in $\ker\phi_0$; and (iii) the two quotient
norm inequalities have the correct directions.  All three are explicit in
the proof.

A bounded novelty search through 11 August 2026 checked the run's registry,
solution, attempt, and proof-gap indexes; arXiv:2507.13866v2; exact-title,
exact-question, author, ideal-null, relative-complementation, and
infinite-quotient searches.  The current v2, revised 18 March 2026, still
contains the question.  No later paper or independent answer was located.
Novelty confidence is moderate, because bounded search is not a proof of
priority.

\section{Human-review recommendation}

Recommended for expert review as a full solution.  The proof is short and
self-contained; review should focus on whether the authors intended any
unstated regularity restriction on the ideals.  Under their printed
definition of proper ideals containing $\Fin$, the construction meets the
question exactly.

\small
\begin{thebibliography}{9}
\bibitem{RU}
M.~A. Rinc\'on-Villamizar and C.~Uzc\'ategui Aylwin,
\emph{On the complementation of spaces of $\mathcal I$-null sequences},
arXiv:2507.13866v2, 2026; Remark~6.7, p.~17.
\end{thebibliography}

\end{document}
